% ---------------------------------------------------------------------------
% How to save an ESM3 structure as a .pdb file (monomer and G-linked trimer)
% and compare it with an experimental structure (2HMG): RMSD, pLDDT, viewers.
%
% Companion to softmax/monte_carlo/structure_comparison_guide.md, which has the
% same content in shorter form. Scripts: decode_to_pdb.py, compare_structures.py.
%
% Build:  cd softmax/docs && pdflatex structure_comparison.tex && pdflatex structure_comparison.tex
% ---------------------------------------------------------------------------
\documentclass[11pt,a4paper]{article}

\usepackage[T1]{fontenc}
\usepackage[utf8]{inputenc}
\usepackage{lmodern}
\usepackage[margin=2.5cm]{geometry}
\usepackage{amsmath}
\usepackage{booktabs}
\usepackage{array}
\usepackage{xcolor}
\usepackage{listings}
\usepackage[hidelinks]{hyperref}
\sloppy

\newcommand{\file}[1]{\texttt{#1}}
\definecolor{codebg}{gray}{0.95}
\lstset{basicstyle=\ttfamily\small, backgroundcolor=\color{codebg}, frame=single,
        rulecolor=\color{gray!40}, columns=fullflexible, keepspaces=true,
        breaklines=true, breakatwhitespace=false, showstringspaces=false,
        upquote=true, xleftmargin=2pt, xrightmargin=2pt}

\title{Saving ESM3 structures as PDB files\\
       and comparing them with 2HMG\\
       \large Monomer and glycine-linked trimer: commands, numbers, viewers}
\author{}
\date{4 November 2026}

\begin{document}
\maketitle

\section{Short answers}

\paragraph{How do I save the structure ESM3 gives me?}
With \file{write\_pdb()} from \file{softmax/monte\_carlo/decode\_to\_pdb.py}.
It is our own helper, not an ESM function. The decode used everywhere in this
repository, \file{decode\_structure\_with\_mask()} in
\file{compute\_structural\_metrics.py}, returns three tensors:
\begin{itemize}
  \item \file{positions}, shape $(L, 37, 3)$: the coordinates of up to 37
        atom types per residue (the ``atom37'' layout), in \AA;
  \item \file{mask}, shape $(L, 37)$: which of those atoms exist;
  \item \file{plddt}, shape $(L)$: the confidence per residue, on 0--1.
\end{itemize}
\file{write\_pdb()} turns these into a text \file{.pdb} file and puts
$100\times$pLDDT in the B-factor column, where every viewer expects it.
You normally do not call it yourself: the script \file{decode\_to\_pdb.py}
decodes and saves in one command (Section~\ref{sec:run}). Section~\ref{sec:python}
shows the call from your own Python code.

\paragraph{Is the trimer with glycine linkers in \file{references.py}?}
No. \file{references.py} holds only single sequences (\file{zero\_polymer},
\file{zero\_polymer\_mature}, the controls). \file{decode\_trimer.py} builds
the trimer on the fly from \file{zero\_polymer}:
\begin{lstlisting}[language=Python]
unit   = ref_seq[16:529]               # residues 17..529, 513 aa
trimer = ("G" * 25).join([unit] * 3)  # 3 x 513 + 2 x 25 = 1589 residues
\end{lstlisting}
\file{decode\_to\_pdb.py} now does the same with \file{-{}-copies 3 -{}-linker 25}.
Note that \file{decode\_trimer.py} takes a \emph{Python slice}
(\file{-{}-residues 16:529}, 0-based, end excluded) while \file{decode\_to\_pdb.py}
takes \emph{residue numbers} (\file{-{}-residues 17:529}, 1-based, both ends
included). Both mean the same 513 residues.

The trimer file is written as three chains A, B, C, one per HA copy, each
numbered 17--529 like the monomer. The glycine linkers are decoded (ESM3 sees
the full 1589-residue sequence) but are \emph{not} written to the file, since
they are not part of HA; their mean pLDDT is printed in the log instead.

\section{What the numbers mean}

\begin{description}
  \item[RMSD] (root-mean-square deviation; this is the ``wmd'' your supervisor
    meant). Pair each residue's CA atom in the prediction with the same
    residue's CA in the experiment, rotate and shift the prediction onto the
    experiment as well as possible (the Kabsch algorithm finds the best rotation
    exactly), then
    \[ \mathrm{RMSD} = \sqrt{\tfrac{1}{N}\textstyle\sum_{i=1}^{N} \lVert x_i^{\mathrm{pred}} - x_i^{\mathrm{exp}} \rVert^2 }. \]
    It is a typical distance in \AA. About 1~\AA: the same structure up to
    noise; 2--3~\AA: same fold, some loops or a domain off; 4--8~\AA: large parts
    placed differently; above about 10~\AA\ for 500 residues: not the same shape.
    One misplaced region inflates it for the whole protein.
  \item[Core RMSD] Fit, drop residues more than 2~\AA\ off, refit, five times
    (what PyMOL's \file{align} reports by default). How well the matching part
    matches; always read it with the fraction of residues it kept.
  \item[Within 1/2/4/8~\AA] Fraction of residues within that distance after the
    global fit. ``70\% within 4~\AA'' is easier to read than an RMSD.
  \item[TM-score] 0 to 1, independent of length. Above 0.5: same fold. Below
    about 0.17: like two unrelated structures.
  \item[Per-chain RMSD] Each experimental chain (HA1, HA2, and for the trimer
    each copy) fitted on its own. Low per chain but high overall means the pieces
    are right but packed wrongly against each other.
  \item[pLDDT] ESM3's confidence per residue: $>90$ very confident, 70--90
    confident, 50--70 low, $<50$ roughly a guess. The script also gives
    Spearman(pLDDT, deviation); negative means the model is least sure exactly
    where it is wrong.
\end{description}

What to expect: HA decodes at mean pLDDT about 0.4--0.6 both as a monomer and as
the G25 trimer (DEVLOG, 2026-11-04), so expect a large RMSD and a TM-score below
0.5. The comparison turns ``the decoder does not fold HA'' into a number against
the real structure. Maltose-binding protein (Section~\ref{sec:control}) shows
what a good match looks like.

\section{Running it on astaroth}\label{sec:run}

All commands are run from \file{softmax/monte\_carlo}; adjust the first
\file{cd} to your clone. Pull first so the scripts are there.

\subsection{The experimental structure}
\begin{lstlisting}[language=bash]
cd ~/my_extendedSampling && git pull
cd softmax/monte_carlo
mkdir -p structure_compare
wget -O structure_compare/2HMG.pdb https://files.rcsb.org/download/2HMG.pdb
grep -c '^ATOM' structure_compare/2HMG.pdb   # thousands of lines = download worked
\end{lstlisting}
2HMG is the cleaved protein: separate HA1 and HA2 chains, three of each.
The comparison script lists the chains it finds and which it used.

\subsection{Predict and save (GPU)}
Monomer (full HA0, the sequence the sampler works on, numbered 1--566):
\begin{lstlisting}[language=bash]
nohup python -u decode_to_pdb.py --ref zero_polymer --num-steps 32 \
    --out structure_compare/ha0_pred.pdb > structure_compare/decode_ha0.log 2>&1 &
\end{lstlisting}
Trimer, exactly as \file{decode\_trimer.py} builds it (ectodomain 17--529,
three copies, 25 glycines between them):
\begin{lstlisting}[language=bash]
nohup python -u decode_to_pdb.py --ref zero_polymer --residues 17:529 \
    --copies 3 --linker 25 --num-steps 32 \
    --out structure_compare/ha_trimer_pred.pdb \
    > structure_compare/decode_trimer.log 2>&1 &
\end{lstlisting}
Follow a log with \file{tail -f structure\_compare/decode\_trimer.log}
(Ctrl-C stops watching, not the job). Each ends with \file{Saved -> ...}, after
printing the mean pLDDT and radius of gyration (per chain and for the linkers,
for the trimer). Run the two decodes one after the other if GPU memory is tight; the
trimer is 1589 residues, the longest decode so far. \file{-{}-num-steps} is the
same setting as in the other decode scripts; HA's structure changes with it, so
write the step count next to every RMSD you report.

\subsection{Compare (CPU, seconds)}
\begin{lstlisting}[language=bash]
python compare_structures.py --pred structure_compare/ha0_pred.pdb \
    --exp structure_compare/2HMG.pdb --out-prefix structure_compare/ha0_vs_2HMG \
    | tee structure_compare/ha0_vs_2HMG.txt

python compare_structures.py --pred structure_compare/ha_trimer_pred.pdb \
    --exp structure_compare/2HMG.pdb --out-prefix structure_compare/trimer_vs_2HMG \
    | tee structure_compare/trimer_vs_2HMG.txt
\end{lstlisting}

How the script pairs residues:
\begin{itemize}
  \item By sequence alignment, not by residue number, so the different numbering
        of our sequence and of 2HMG does not matter.
  \item \textbf{Monomer.} ESM3 gives one chain (HA0); 2HMG has HA1 and HA2 as
        separate chains. Both are aligned onto the one predicted chain. The other
        two copies in 2HMG would land on the same predicted residues, so they are
        skipped (``symmetry copy'').
  \item \textbf{Trimer.} Residues already paired are hidden from the next
        alignment, so the three HA1/HA2 pairs of 2HMG land on the three
        predicted chains A, B, C. Which experimental copy goes with which
        predicted copy is arbitrary, so the script tries every order and keeps
        the lowest RMSD; the line ``predicted protomers tried in every order''
        lists all of them. The trimer number therefore includes how the three
        copies are packed, the thing the trimer was built to test.
  \item To compare one predicted copy only: \file{-{}-pred-chain A}. To choose
        experimental chains by hand: \file{-{}-exp-chains A,B}.
  \item Residues absent from the crystal (signal peptide, transmembrane anchor,
        tail, disordered loops) cannot be compared and are left out of every
        number; the script says how many.
\end{itemize}

Files written next to the prefix:
\begin{center}
\begin{tabular}{>{\ttfamily}l p{0.55\textwidth}}
\toprule
\_superposed.pdb   & the prediction moved on top of 2HMG \\
\_per\_residue.csv & per residue: predicted chain and number, pLDDT, matching 2HMG chain and number, deviation \\
\_per\_residue.png & deviation (blue) and pLDDT (red) along the file; dashed lines separate experimental chains \\
\_view.cxc / .pml  & ready-made ChimeraX / PyMOL scenes \\
\bottomrule
\end{tabular}
\end{center}

\subsection{Positive control}\label{sec:control}
Maltose-binding protein decodes at pLDDT 0.93; PDB 1ANF is its crystal structure.
\begin{lstlisting}[language=bash]
wget -O structure_compare/1ANF.pdb https://files.rcsb.org/download/1ANF.pdb
nohup python -u decode_to_pdb.py --uniprot P0AEX9 --num-steps 32 \
    --out structure_compare/mbp_pred.pdb > structure_compare/decode_mbp.log 2>&1 &
python compare_structures.py --pred structure_compare/mbp_pred.pdb \
    --exp structure_compare/1ANF.pdb --out-prefix structure_compare/mbp_vs_1ANF
\end{lstlisting}
\file{-{}-uniprot} needs network access to UniProt and takes the mature chain.

\section{Calling it from your own Python code}\label{sec:python}

If you already have the tensors from a decode inside your own script (for
example a sampled sequence), save them like this. Run from
\file{softmax/monte\_carlo} so the imports resolve, as in the other scripts.
\begin{lstlisting}[language=Python]
from compute_structural_metrics import decode_structure_with_mask
from decode_to_pdb import write_pdb
from custom_esm.utils.residue_constants import atom_types

positions, mask, plddt = decode_structure_with_mask(sampler.model, eprot, config)

write_pdb("my_structure.pdb", eprot.sequence,
          positions.float().numpy(), mask.bool().numpy(),
          plddt.float().numpy() * 100,     # pLDDT 0-100 into the B-factor column
          atom_types)
\end{lstlisting}
For a G-linked trimer, pass \file{segments} so each copy becomes its own chain
and the linkers are left out. Each entry is (chain, first index, end index,
first residue number), with 0-based indices into the decoded sequence:
\begin{lstlisting}[language=Python]
n, L = 513, 25                          # residues per copy, glycines per linker
segments = [("A", 0,           n,           17),
            ("B", n + L,       2 * n + L,   17),
            ("C", 2 * (n + L), 3 * n + 2*L, 17)]
write_pdb("trimer.pdb", eprot.sequence,
          positions.float().numpy(), mask.bool().numpy(),
          plddt.float().numpy() * 100, atom_types, segments=segments)
\end{lstlisting}
ESM's own \file{ESMProtein.to\_pdb()} also writes PDB files, but the decode
path here never builds an \file{ESMProtein} with coordinates, and the helper
the repository used before (\file{test\_cleaved\_complex.py -{}-pdb}) leaves
every B-factor at 1.00 and writes pLDDT to a separate \file{.plddt.tsv}.
\file{write\_pdb()} keeps everything in one file.

\section{Looking at the structures}

A viewer draws 3D graphics, so it needs a graphical session on the machine that
runs it. astaroth has a screen, so pick by where you are sitting.

\paragraph{At astaroth's keyboard} install the viewer there and open the files in
place. Nothing to copy.

\paragraph{Over ssh} either render the picture on astaroth and copy back only the
image, which needs no graphics session at all,
\begin{lstlisting}[language=bash]
chimerax --offscreen --exit --cmd "open trimer_vs_2HMG_view.cxc ; \
    save view.png width 2000 supersample 3"
\end{lstlisting}
(\texttt{--offscreen} needs a build with OSMesa, which the Linux builds normally
have; if yours does not it says so immediately), or copy the folder to your
laptop and look at it there (run this from the laptop; replace the login and
path):
\begin{lstlisting}[language=bash]
scp -r marco.sanchini@astaroth:~/my_project/extendedSampling/softmax/monte_carlo/structure_compare .
cd structure_compare
\end{lstlisting}
Copy the whole folder, not single files: the \file{.cxc} and \file{.pml} scripts
refer to the others by relative path. \texttt{ssh -X} also works but forwards
every frame, so it is only good for checking that a file opens.

\paragraph{ChimeraX} (free for academic use,
\url{https://www.cgl.ucsf.edu/chimerax/download.html}). Quickest: \emph{File
$>$ Open\ldots} and pick \file{trimer\_vs\_2HMG\_view.cxc}. 2HMG appears in
grey, the prediction coloured by pLDDT (dark blue $>90$, light blue 70--90,
yellow 50--70, orange $<50$), already on top of each other. By hand, in the
command line at the bottom of the window:
\begin{lstlisting}
open 2HMG.pdb                       # model #1
open ha_trimer_pred.pdb             # model #2
delete solvent
matchmaker #2/A to #1/A             # superimpose; prints an RMSD in the log
color #1 lightgray
color bfactor #2 palette alphafold
view
\end{lstlisting}
\file{matchmaker} fits one chain to one chain (here predicted copy A onto HA1
chain A), so its RMSD differs from the script's, which fits everything at once.
It prints two values: over ``pruned'' pairs (like the core RMSD) and ``across all
pairs''. \file{save picture.png width 2000 supersample 3} saves an image.

\paragraph{PyMOL} (\file{conda install -c conda-forge pymol-open-source}).
\emph{File $>$ Run Script\ldots} and pick \file{trimer\_vs\_2HMG\_view.pml}, or
type:
\begin{lstlisting}
load 2HMG.pdb, exp
load ha_trimer_pred.pdb, pred
remove solvent
align pred and name CA, exp and name CA, cycles=0   # all residues, no outlier rejection
hide everything
show cartoon
color grey70, exp
spectrum b, red_white_blue, pred, minimum=30, maximum=90
orient
\end{lstlisting}
The log line \file{Executive: RMSD = \ldots} gives the RMSD.

\paragraph{Browser, nothing to install.} \url{https://molstar.org/viewer/}: drag
in \file{2HMG.pdb} and the \file{\_superposed.pdb} file (Mol* does not fit them
for you), then set the prediction's colour theme to ``Uncertainty/Disorder''
to colour by the B-factor column.

\section{Things that trip people up}
\begin{itemize}
  \item \textbf{Slice versus residue numbers}: \file{16:529} in
        \file{decode\_trimer.py} is \file{17:529} in \file{decode\_to\_pdb.py}.
  \item \textbf{One chain against two}: HA0 is one chain in the prediction, HA1 +
        HA2 in the crystal. The script handles it; single-chain tools such as
        ChimeraX's \file{matchmaker} give different numbers.
  \item \textbf{Monomer against trimer}: a monomer prediction has no partner
        copies holding the stalk helices, so part of its disagreement with 2HMG is
        quaternary. Compare the per-chain RMSDs of the monomer and trimer runs.
  \item \textbf{Step count}: report \file{-{}-num-steps} with every number.
  \item \textbf{Strain}: the identity line says how close 2HMG's sequence is to
        ours; mention it if it is well below 100\%.
\end{itemize}

\section{The neuraminidase candidate}\label{sec:na}

HA0 failed the decode gate, so the sampling target is being chosen again
(\file{target\_selection.md}). The first candidate is the influenza neuraminidase
catalytic head, and unlike HA it comes with an experimental structure of the
construct itself: 2HU0 chain A, supplied by the supervisor stripped to one chain,
kept at \file{code/bridge\_test/pdb/na\_head\_2HU0\_monomer.pdb}. Its sequence is
\file{references.py} key \file{na\_head\_2hu0}, 385\,aa.

Everything that made HA awkward is absent here: one predicted chain against one
experimental chain, no cleavage, no trimer, and no waters or ligands left in the
file, so nothing has to be stripped and \file{matchmaker \#2 to \#1} needs no
chain selector.
\begin{lstlisting}[language=bash]
python decode_to_pdb.py --ref na_head_2hu0 --num-steps 32 \
    --out structure_compare/na_head_pred.pdb

python compare_structures.py \
    --pred structure_compare/na_head_pred.pdb \
    --exp ../code/bridge_test/pdb/na_head_2HU0_monomer.pdb \
    --out-prefix structure_compare/na_head_vs_2HU0
\end{lstlisting}

How to read it. A mean pLDDT near $0.9$ with a low RMSD means the candidate
passes the gate and sampling can start on it; a pLDDT near HA0's $0.5$ means drop
it and move to the next candidate. The picture carries information the numbers do
not: a six-bladed $\beta$-propeller is a pinwheel, so rotate until you look down
the central axis and six four-stranded sheets should fan out around it. A
prediction that is compact but shows no pinwheel has failed even if its RMSD is
not catastrophic. Run \file{-{}-num-steps 8} and \file{128} as well, to different
\file{-{}-out} names: a confident decode barely moves between step counts, while
HA0 wandered from $R_g$ 45 to 70\,\AA.

Traps specific to this pair, on top of the list above:
\begin{itemize}
  \item \textbf{Numbering}: the crystal runs 83--468 with five gaps totalling six
        unobserved loop residues and insertion codes 169A and 412A--D, while the
        prediction runs 1--385. The offset is not a constant 82 and changes at
        every gap, so take the pairing from
        \file{na\_head\_vs\_2HU0\_per\_residue.csv} and never compute it by hand.
  \item \textbf{The reference is the observed sequence}: those six residues are
        missing from \file{na\_head\_2hu0}, which closes the gaps up. Acceptable
        for the gate, since all six are in loops, but before any sampling run or
        K-site comparison replace it from 2HU0's SEQRES
        (\file{pdb\_to\_reference.py -{}-use-seqres} on the unstripped file) or
        from the strain's UniProt entry (\file{add\_reference.py}).
  \item \textbf{A tetramer, but not HA's problem}: native NA is a homotetramer
        with a Ca$^{2+}$ site on the fourfold axis, so a monomer prediction is
        missing partners here too. The difference is that the head's fold is
        entirely intra-chain, a self-contained $\beta$-propeller, whereas HA2's
        long helix exists only as part of the trimeric coiled coil. Expect the
        interface surface to be the loosest part of the prediction, and flag
        interface residues before the per-site entropy comparison: they are
        conserved for quaternary reasons a single-chain energy cannot see.
\end{itemize}

\end{document}
